% BHU PYQ -- Manually transcribed & error-corrected
% Paper: GLB-503: Sedimentary and Metamorphic Petrology
% Exam: B.Sc. (Hons.) Semester V Examination, 2013-14

\documentclass[11pt,a4paper]{article}
\usepackage[a4paper,top=2.5cm,bottom=2.5cm,left=2.8cm,right=2.8cm,headheight=14pt]{geometry}
\usepackage{amsmath,amssymb}
\usepackage[T1]{fontenc}
\usepackage[utf8]{inputenc}
\usepackage{lmodern,microtype,fancyhdr,enumitem,hyperref,lastpage,parskip}

\hypersetup{colorlinks=true,urlcolor=black,linkcolor=black,
  pdftitle={GLB-503 Sedimentary and Metamorphic Petrology -- BHU B.Sc. Sem V 2013-14}}
\pagestyle{fancy}\fancyhf{}
\renewcommand{\headrulewidth}{0.4pt}\renewcommand{\footrulewidth}{0.4pt}
\fancyhead[L]{\small   \textit{Banaras Hindu University}}
\fancyhead[R]{\small   \textit{GLB-503: Sedimentary and Metamorphic Petrology}}
\fancyfoot[C]{\small Page  \thepage\ of \pageref{LastPage}}


\newlist{parts}{enumerate}{2}
\setlist[parts,1]{label=(\alph*),leftmargin=*,itemsep=5pt,topsep=3pt}

\newcommand{\pts}[1]{\hfill   \textit{[#1]}}


\begin{document}

\begin{center}
  {\large\bfseries BANARAS HINDU UNIVERSITY}\\[2pt]
  {
\normalsize B.Sc. (Hons.) --- Semester V Examination, 2013-14}\\[6pt]
  \rule{0.8\linewidth}{0.5pt}\\[6pt]
  {\large\bfseries Paper No. GLB-503: Sedimentary and Metamorphic Petrology}\\[4pt]
  {
\normalsize Time: 3 Hours \hfill Maximum Marks: 70}\\[2pt]
  {\small Ref. No.: BS/Sem V/259}
\end{center}
\rule{\linewidth}{0.4pt}
\smallskip



\noindent   \textit{Instructions: Attempt any   \textbf{five} questions, selecting at least   \textbf{two} questions each from Section A and Section B. Section C is compulsory. All questions carry equal marks (14 marks each).}
\bigskip

\bigskip


\noindent {\large\bfseries SECTION --- A}
\rule{\linewidth}{0.4pt}\smallskip




\noindent   \textbf{Question 1.} Describe the process of diagenesis in sandstones. \pts{14}

\medskip


\noindent   \textbf{Question 2.} With suitable sketches, describe the internal sedimentary structures. \pts{14}

\medskip


\noindent   \textbf{Question 3.} Write notes on   \textbf{any two} of the following: \pts{$2  \times 7 = 14$}

\begin{parts}
  \item Ripple marks
  \item Folk's classification of limestones
  \item Sediment deposition
\end{parts}

\bigskip


\noindent {\large\bfseries SECTION --- B}
\rule{\linewidth}{0.4pt}\smallskip




\noindent   \textbf{Question 4.} Describe the progressive metamorphism of Pelitic rocks in the $ \text{K}_2 \text{O}- \text{FeO}- \text{MgO}- \text{Al}_2 \text{O}_3- \text{SiO}_2$ system. \pts{14}

\medskip


\noindent   \textbf{Question 5.} Discuss in detail the metamorphic facies. \pts{14}

\medskip


\noindent   \textbf{Question 6.} Write notes on   \textbf{any two} of the following: \pts{$2  \times 7 = 14$}

\begin{parts}
  \item Retrograde Metamorphism
  \item Phase rule
  \item Principles of Metamorphic reactions
\end{parts}

\medskip


\noindent   \textbf{Question 7.} Describe progressive metamorphism of Basic rocks in the $ \text{CaO}- \text{FeO}- \text{MgO}- \text{Al}_2 \text{O}_3- \text{SiO}_2$ system. \pts{14}

\bigskip


\noindent {\large\bfseries SECTION --- C (Compulsory)}
\rule{\linewidth}{0.4pt}\smallskip




\noindent   \textbf{Question 8.} Fill in the blanks or mark True (T) and False (F) as indicated: \pts{$7  \times 2 = 14$}

\begin{parts}
  \item The typical product of contact metamorphism with maculose structure is \underline{\hspace{4cm}}.
  \item Hot and humid climate promote physical weathering (T/F). \underline{\hspace{1cm}}
  \item Charnockite is a typical rock of \underline{\hspace{3cm}} facies.
  \item Graywacke sandstones are mineralogically and texturally \underline{\hspace{4cm}}.
  \item Foliation is defined by the \underline{\hspace{4cm}} arrangement of minerals.
  \item Flute casts are good palaeocurrent indicators (T/F). \underline{\hspace{1.5cm}}
  \item Greenschist facies represents \underline{\hspace{3cm}} grade metamorphism.
\end{parts}

\vfill\rule{\linewidth}{0.4pt}

\begin{center}\small
\href{https://bhu-exam.vercel.app/}{Click here to visit our website}\\[4pt]
\href{https://me-aryan.vercel.app/}{Click here to visit our contributor}\\[4pt]
\href{https://quashan-me.vercel.app/}{Click here to visit our contributor}
\end{center}
\end{document}
